% \section{Related Work}

% \lijie{We need to shorten these two paragraphs, only compare the close related work.}


% \lijie{Just focus on the RL-based methods and mention other methods are RL-free and orthogonal.} \yiran{RL-free method also aims to improve model reasoning capacity. So they are not orthogonal? (Two categories of methods, can't use together) 
% We also compare RestEM in the SPO-chain experiment.} 









% For instance, RFT \cite{yuan2024advancingllmreasoninggeneralists, yuan2023scalingrelationshiplearningmathematical} fine-tunes the pretrained base model using correct samples generated by a supervised fine-tuned model. Similarly, RestEM \cite{singh2024humandatascalingselftraining} adopts an iterative self-training strategy, repeatedly retraining the base model on high-quality reasoning traces produced by previously obtained checkpoints. Preference-based methods like DPO \cite{rafailov2024directpreferenceoptimizationlanguage,rafailov2024rqlanguagemodel} optimize models by contrasting correct and incorrect reasoning outputs. Search-guided approaches, like Monte Carlo Tree Search (MCTS) \cite{zhang2024restmctsllmselftrainingprocess, chen2024alphamathzeroprocesssupervision}, help models discover improved reasoning paths during inference. Another prominent direction involves reinforcement learning (RL) frameworks. These frameworks differ mainly in advantage estimation granularity. For example, GRPO \cite{shao2024deepseekmathpushinglimitsmathematical}, RLOO \cite{ahmadian2024basicsrevisitingreinforcestyle} estimate trajectory-level advantages, while PPO \cite{schulman2017proximalpolicyoptimizationalgorithms, ouyang2022traininglanguagemodelsfollow} estimates token-level advantages using a dedicated critic model. Recent work \cite{kazemnejad2024vineppounlockingrlpotential} highlights that accurately training the critic in PPO is challenging, and proposes VinePPO, a critic-free alternative that partitions the reasoning trajectory into discrete steps using heuristic rules (e.g., line breaks) and estimates step-level advantages through Monte Carlo (MC) sampling. Compared to VinePPO, our segment-level framework adopts a more general concept of segmentation, allowing arbitrary partitions without enforcing semantic coherence. This enables us to freely adjust the granularity anywhere between token-level and trajectory-level, also facilitating adaptation to broader and less structured tasks such as code generation. Notably, VinePPO can be viewed as a special case within our more general SPO framework.





% Compared to VinePPO, our segment-level framework adopts a broader and more general concept of segmentation. Rather than relying on manually crafted heuristic rules or semantic coherence, our method allows arbitrary segment partitions, producing segments that can be either finer or coarser in granularity than predefined reasoning steps. This generality enables our method to be applicable beyond scenarios where reasoning trajectories can be clearly divided into semantic steps, extending naturally to broader and less structured tasks such as code generation. Moreover, we introduce a novel tree-based segment advantage estimation method within our unified framework, significantly reducing the computational costs associated with MC-based advantage computation. This effectively resolves the key challenge of expensive MC estimation, allowing our method to scale effortlessly to long CoT scenario. Notably, VinePPO can be considered as a special case instantiated within our more general SPO framework.














% Crucially, these studies focus on training an external reward model (Process Reward Model, or PRM), which serves two key purposes: providing intermediate rewards to guide policy optimization and facilitating Best-of-N (BoN) selection at inference time. In practice, when applying PRM for policy optimization, existing approaches typically combine these step-level rewards with the final binary correctness reward and still employ standard RL algorithms such as PPO or GRPO for training. In contrast, our approach differs fundamentally by improving the RL optimization algorithm itself (e.g., PPO/GRPO), rather than relying on external reward models. Specifically, we introduce a Monte Carlo (MC) based estimation method to directly compute segment-level advantages from the current policy, aiming at resolving gradient estimation variance and enabling finer-grained credit assignment during optimization. This strategy helps stabilize training and leads to more effective policy updates compared to using sparse, trajectory-level rewards alone. Notably, our MC advantage estimation framework remains fully compatible with the PRM approach: intermediate process reward signals generated via PRMs could be integrated in our trajectory evaluations. Such integration may combine the strengths of both methodologies, providing further potential improvements in overall model performance. In Appendix \ref{appendix:incorporate_process_reward}, we provide a specific method for integrating Process Reward into our framework.